\section{Praktikum}
\label{sec:praktikum-bab8}

\begin{praktikumbox}
\textbf{Judul:} Penelusuran mode pengalamatan dan verifikasi alamat
efektif dengan \eng{debugger}.\\[0.3em]
\textbf{Sub-CPMK:} 2.2 \quad\textbullet\quad
\textbf{Alokasi waktu:} 100 menit \quad\textbullet\quad
\textbf{Bentuk:} individu
\end{praktikumbox}

\subsection{Tujuan Praktikum}

Setelah menyelesaikan praktikum ini, mahasiswa mampu menelusuri
eksekusi program satu instruksi demi satu, menghitung alamat efektif
setiap operand memori, dan memverifikasinya melalui jendela memori.

\subsection{Alat dan Bahan}

\begin{itemize}
    \item Komputer dengan sistem operasi Windows.
    \item Emu8086 (emulator dengan tampilan register dan memori).
    \item DOSBox beserta \file{DEBUG} untuk penelusuran berbasis
          teks, serta TASM dan TLINK untuk merakit program.
    \item Berkas \file{mode-array.asm} dan \file{param-bp.asm} dari
          \file{code/asm/}.
    \item Lembar kerja pengamatan alamat efektif.
\end{itemize}

\subsection{Langkah Kerja}

\begin{enumerate}
    \item \textbf{Menyiapkan program.} Rakit kedua berkas contoh
          dengan TASM dan TLINK, lalu pastikan keluarannya sesuai
          Kode~\ref{lst:keluaran-mode-array} dan
          Kode~\ref{lst:keluaran-param-bp}.

    \item \textbf{Menelusuri dengan Emu8086.}
    \begin{enumerate}
        \item Buka \file{mode-array.asm}, tekan \file{Compile and
              Emulate}, lalu aktifkan tombol \file{Single Step}.
        \item Catat nilai \reg{DS}, \reg{BX}, \reg{SI}, dan \reg{BP}
              sebelum setiap instruksi memori dieksekusi, lalu hitung
              alamat efektifnya secara manual.
        \item Bandingkan hasil hitungan Anda dengan alamat yang
              ditampilkan emulator.
    \end{enumerate}

    \item \textbf{Menelusuri dengan \file{DEBUG}.}
    \begin{enumerate}
        \item Jalankan DOSBox, lalu ketik
              \file{debug mode-array.exe}.
        \item Tekan \file{u} untuk melihat \eng{disassembly} dan
              \file{r} untuk melihat isi register.
        \item Tekan \file{t} beberapa kali dan catat perubahan
              register setiap instruksi.
        \item Tampilkan \file{d ds:0} dan \file{d ss:0}, lalu
              verifikasi alamat yang ditunjuk benar-benar berisi nilai
              array, bukan nol atau karakter sampah.
    \end{enumerate}

    \item \textbf{Membandingkan \code{[BX]} dan \code{[BP]}.}
    \begin{enumerate}
        \item Pada \file{param-bp.asm}, jalankan sampai
              \lstasm|MOV BP, SP|, lalu catat nilai \reg{SS} dan
              \reg{SP}.
        \item Ubah \lstasm|MOV AX, [BP+0]| menjadi
              \lstasm|MOV AX, [BX+0]| dengan \reg{BX} bernilai nol.
        \item Rakit ulang dan jelaskan mengapa keluarannya berubah,
              dengan mengaitkannya pada
              Tabel~\ref{tab:segmen-default}.
    \end{enumerate}

    \item \textbf{Mengisi lembar pengamatan.} Lengkapi
          Tabel~\ref{tab:pengamatan-ea} berdasarkan hasil penelusuran,
          lalu kerjakan Kuis 2 pada akhir praktikum.
\end{enumerate}

\begin{table}[H]
    \centering
    \caption{Lembar pengamatan alamat efektif yang harus diisi
    mahasiswa.}
    \label{tab:pengamatan-ea}
    \small
    \begin{tabular}{@{}p{0.7cm}p{3.7cm}p{2.7cm}p{1.7cm}p{2.2cm}p{1.8cm}@{}}
        \toprule
        \textbf{No} & \textbf{Instruksi} & \textbf{EA terhitung} &
        \textbf{Segmen} & \textbf{Isi memori} &
        \textbf{Sesuai} \\
        \midrule
        1 & \lstasm|MOV AX, data1+0| & \code{offset data1} & \reg{DS}
        & \hex{000A} (10) & Ya \\
        \addlinespace
        2 & \lstasm|MOV AX, data1[SI]|, \reg{SI} = 2 &
        \dots & \dots & \dots & \dots \\
        \addlinespace
        3 & \lstasm|MOV AX, [BX+SI]|, \reg{SI} = 4 &
        \dots & \dots & \dots & \dots \\
        \addlinespace
        4 & \lstasm|MOV AX, [BP+0]| pada \file{param-bp.asm} &
        \dots & \dots & \dots & \dots \\
        \addlinespace
        5 & \lstasm|MOV AX, [BP+2]| pada \file{param-bp.asm} &
        \dots & \dots & \dots & \dots \\
        \bottomrule
    \end{tabular}
\end{table}

\subsection{Hasil yang Dikumpulkan}

\begin{enumerate}
    \item Tabel~\ref{tab:pengamatan-ea} yang telah dilengkapi.
    \item Tangkapan layar jendela register dan memori Emu8086 untuk
          minimal tiga instruksi berbeda.
    \item Tangkapan layar sesi \file{DEBUG} dengan perintah
          \file{u}, \file{t}, dan \file{d}, serta lembar jawaban
          Kuis 2.
\end{enumerate}

\subsection{Kuis 2}

Kuis 2 dilaksanakan pada 15 menit terakhir praktikum secara individu
tanpa membuka catatan.

\begin{enumerate}
    \item Tentukan mode pengalamatan dari lima penulisan berikut:
          \lstasm|MOV AX, 20|, \lstasm|MOV AX, BX|,
          \lstasm|MOV AX, [BX]|, \lstasm|MOV AX, [SI+4]|, dan
          \lstasm|MOV AX, [BX+DI+2]|.

    \item Diketahui \reg{DS} = \hex{2000}, \reg{BX} = \hex{0100},
          \reg{SI} = \hex{0004}, dan \reg{SS} = \hex{3000}.
          Hitung alamat efektif dan alamat fisik untuk
          \lstasm|MOV AX, [BX+SI+8]| serta
          \lstasm|MOV AX, [BP]| dengan \reg{BP} = \hex{0010}.

    \item Program membaca variabel segmen data dengan
          \lstasm|MOV AX, [BP+2]| dan hasilnya salah. Jelaskan
          penyebabnya dan tuliskan dua cara memperbaikinya tanpa
          mengubah nilai perpindahannya.
\end{enumerate}

\begin{asesmenbox}
\textbf{Asesmen praktikum 2 dan Kuis 2 (bobot 10\% dari komponen
praktikum).}

\begin{tabular}{@{}p{5.9cm}cp{4.9cm}@{}}
    \toprule
    \textbf{Indikator} & \textbf{Bobot} & \textbf{Kriteria} \\
    \midrule
    Ketepatan perhitungan alamat efektif & 25\% &
    Seluruh baris lembar pengamatan terisi benar \\
    \addlinespace
    Verifikasi melalui jendela memori & 20\% &
    Alamat dan isi memori yang dicatat sesuai dengan
    \eng{debugger} \\
    \addlinespace
    Analisis perbedaan \code{[BX]} dan \code{[BP]} & 15\% &
    Penjelasan menyebut peran segmen \reg{DS} dan \reg{SS} \\
    \addlinespace
    Kuis 2 & 30\% &
    Minimal tiga dari empat bagian soal terjawab benar \\
    \addlinespace
    Kerapian dokumentasi praktikum & 10\% &
    Tangkapan layar lengkap dan langkah kerja dapat diulang \\
    \bottomrule
\end{tabular}
\end{asesmenbox}
